% Proposed NS024 appendix insertion; reviewed integration and compilation are separate.
\section{LLM use, public reproduction, and responsible use}
\label{sec:ns024-disclosures}

\paragraph{Scientific model use.}
The final cold-repair comparison uses Grok 4.6 to propose source edits from
the fixed A/B context requests (Section~\ref{sec:design}). Each of the 32 cells
consumes one proposal POST with a maximum of 4,096 requested output tokens and
a structured edit response. No scientific fallback, retry, served random seed,
or reasoning-effort parameter is used. The model's weights and exact served
revision are unavailable. Candidate inspection by AI is a procedural control;
it is not independent human assessment or a substitute for executed checks.

\paragraph{Research and manuscript assistance.}
AI assistance also contributed to source recovery and drafting, implementation
and debugging, experiment preparation and candidate inspection, and analysis
and artifact preparation. The authoring supervisors were configured to use
Grok 4.6, with Codex \texttt{gpt-5.6-terra} at high reasoning effort as a fallback
only following independently verified primary-quota exhaustion. This describes
the configured authoring policy, not a claim that every task used a fallback or
that these models generated the scientific comparison under a mixed route.
The configured policy does not identify the model behind every separate
interactive assistance session. AI-assisted checking supplies no measured
human semantic fidelity, expert agreement, or human effort time, and does not
discharge the conditional arguments in Appendix~\ref{sec:formal-arguments-ns023}.

\paragraph{Anonymous numerical reproduction.}
The accompanying supplement keeps the 32-cell final comparison, the separate
24-cell developmental pilot, and historical boundary controls in distinct
components. Its top-level \texttt{README.md} specifies Python 3.12 and the
command \texttt{python3 -B reproduce.py --root ROOT --output OUTPUT}, where
\texttt{ROOT} is a fresh extraction and \texttt{OUTPUT} is a fresh directory
outside it. The program checks the manifest-bound files and reproduces their
retained numerical summaries; it does not generate new proposals, execute
candidates or hidden tests, rerun historical qualifications, or authenticate
original signed records. Hidden tests, reference patches, credentials, keys,
and raw private operational logs are excluded. Original and derivative hashes
have distinct roles: original signatures do not authenticate modified anonymous
derivative bytes. Exact executed source and labeled reference derivatives remain
distinct. Upstream retrieval metadata, revisions, and retained license and
copyright notices accompany the released assets. Stable hashes and timestamps
may permit correlation, so these release controls do not guarantee anonymity.

\paragraph{Potential impacts and limits.}
Explicit evidence boundaries may help users identify stale checks and
unsupported acceptance claims, but this experiment establishes no general
reliability or efficiency benefit. Incorrect repairs, incomplete tests, or
mistaking signed provenance for correctness could lead to unsafe deployment;
automated coding can also be misused to introduce harmful changes. The stated
scope, isolated checks, restricted artifact projection, and visible failures
reduce particular risks without establishing comprehensive adversarial
resistance. No model weights or human-subject dataset are released, and no
human semantic annotations, expert-agreement study, or human-time measurements
were collected for the reported comparison. 